% Janmitra final project report. This is the only required source file.
% Compile once with Tectonic (which performs reruns), or twice with pdfLaTeX.
% Structure follows the prototype-stage report. References and the architecture
% figure are embedded; no external assets are needed.
\documentclass[11pt,a4paper,oneside]{report}

\usepackage[T1]{fontenc}
\usepackage[utf8]{inputenc}
\usepackage{lmodern}
\usepackage[margin=1in,headheight=15pt]{geometry}
\usepackage{microtype}
\usepackage{enumitem}
\usepackage{array,booktabs,tabularx,longtable}
\usepackage{xcolor}
\usepackage{amsmath}
\usepackage{fancyhdr}
\usepackage{xurl}
\usepackage[hidelinks]{hyperref}

\hypersetup{
  pdftitle={Janmitra: Final Project Report},
  pdfauthor={Dhruv Srivastava; Harkamal Singh Lubana; Paras},
  pdfsubject={UCS503P final implementation and evaluation}
}
\urlstyle{tt}
\setlength{\parindent}{0pt}
\setlength{\parskip}{0.45em}
\setlength{\emergencystretch}{3em}
\renewcommand{\arraystretch}{1.18}
\setlist[itemize]{leftmargin=1.5em,itemsep=0.2em,topsep=0.3em}
\setlist[enumerate]{leftmargin=1.8em,itemsep=0.25em,topsep=0.3em}
\newcolumntype{Y}{>{\raggedright\arraybackslash}X}
\newcolumntype{P}[1]{>{\raggedright\arraybackslash}p{#1}}
\newcommand{\code}[1]{\texttt{\detokenize{#1}}}
\newcommand{\ReportDate}{27 November 2026}
\newcommand{\Revision}{ac5c166}
\newcommand{\Advisor}{Ms Paramveer Kaur}

\pagestyle{fancy}
\fancyhf{}
\fancyhead[L]{\small Janmitra}
\fancyhead[R]{\small Final Project Report}
\fancyfoot[C]{\thepage}
\renewcommand{\headrulewidth}{0.4pt}
\fancypagestyle{plain}{%
  \fancyhf{}\fancyfoot[C]{\thepage}\renewcommand{\headrulewidth}{0pt}}

\begin{document}

\hypersetup{pageanchor=false}
\begin{titlepage}
  \centering
  {\large\bfseries Thapar Institute of Engineering and Technology\par}
  \vspace{0.2cm}
  {\large Patiala\par}
  \vspace{0.9cm}
  {\large UCS503P Course Project\par}
  \vspace{1.1cm}
  {\Huge\bfseries Janmitra\par}
  \vspace{0.4cm}
  {\LARGE A Voice-First Civic Scheme Guidance\par
    Platform for Rural India\par}
  \vspace{0.8cm}
  {\Large\bfseries Final Project Report\par}
  \vspace{0.2cm}
  {\large End-Semester Evaluation\par}
  \vspace{1.2cm}
  {\large Submitted to\par}
  \vspace{0.15cm}
  {\Large\bfseries \Advisor\par}
  \vspace{1.0cm}
  {\large Submitted by\par}
  \vspace{0.3cm}
  \begin{tabular}{@{}ll@{}}
    \textbf{Student} & \textbf{Roll number} \\
    \midrule
    Dhruv Srivastava & \texttt{1024170394} \\
    Harkamal Singh Lubana & \texttt{1024170396} \\
    Paras & \texttt{1024170395}
  \end{tabular}
  \vspace{0.4cm}
  \par{\small
    \texttt{dsrivastava\_be24@thapar.edu}\par
    \texttt{hlubana\_be24@thapar.edu}\par
    \texttt{pparas\_be24@thapar.edu}\par}
  \vfill
  {\large\bfseries Computer Science and Engineering Department\par}
  \vspace{0.4cm}
  {\large \ReportDate\par}
\end{titlepage}

\pagenumbering{roman}
\hypersetup{pageanchor=true}
\chapter*{Abstract}
\addcontentsline{toc}{chapter}{Abstract}
Janmitra is a voice-first civic-guidance system that helps rural citizens discover
government support, understand eligibility conditions, prepare documents, and
request human assistance. This report describes the final implementation at
repository revision \texttt{\Revision}. Citizens arrive on a scroll-driven
landing page that maps the 18 published schemes onto concentric rings, then speak
through a phone-call screen connected over LiveKit, where Gemini Live handles
native speech input and output. A FastAPI backend provides scheme retrieval, deterministic eligibility
evaluation, document checklists, versioned publication, an operator handoff queue,
and conversation and audit records.

The final evaluation covered automated testing, end-to-end API verification,
catalogue coverage, retrieval accuracy, load and replica scaling, and live voice
trials. All backend tests passed (100\%), including PostgreSQL integration and
voice-runtime modules, as did all frontend tests (100\%), including new tests of
the landing page, and every continuous-integration workflow (100\%). End-to-end
API checks passed at 96.9\%. All 18 catalogue records were published with a named reviewer
attestation, and every official source link resolved (100\%). Retrieval reached
97.1\% recall at three results, including 100\% on Hindi, Punjabi, and
transliterated queries. Throughput exceeded each replica-scaling target by more
than tenfold with a 0\% error rate.

In live voice trials, backend retrieval succeeded in every call (100\%), and
50\% of integrated calls completed; the remainder were closed by the speech
provider's free-tier capacity limit. A Punjabi call demonstrated correct language
matching and cross-language retrieval. Telephone access, a median turn-level
latency measurement, and a participant pilot remain future work.

\tableofcontents
\clearpage
\pagenumbering{arabic}

% ================================================================
\chapter{Introduction}

\section{Background and Context}
Government welfare information is distributed across departmental websites,
guidelines, notices, and application portals. A citizen may describe a practical
need, such as support for farming, employment, housing, insurance, or a
small business, without knowing the relevant scheme or department. Formal language,
unfamiliar terminology, and multi-page navigation can make public information
difficult to use even when it is available.

Janmitra aims to improve access to this information through conversation. The
original proposal envisaged a service reachable by an ordinary phone call, so that
citizens could speak in a familiar language without installing an application
\cite{proposal}. The project developed the speech and backend workflow through a
browser voice channel presented as a phone-call screen, which provides a practical
environment for integration, measurement, and demonstration. This report extends
the prototype-stage report \cite{prototype} with the final implementation and its
measured evaluation.

\subsection{Problem Statement}
The project addresses five connected difficulties:
\begin{enumerate}
  \item \textbf{Fragmented information:} relevant facts are spread across
    multiple official sources and versions.
  \item \textbf{Language and literacy barriers:} a citizen may find it easier to
    describe a need aloud than to read formal scheme text.
  \item \textbf{Discovery difficulty:} a person may know their circumstances but
    not the name of a suitable programme.
  \item \textbf{Eligibility and document uncertainty:} general descriptions do not
    explain which conditions apply to a particular person or what to prepare.
  \item \textbf{Incomplete assistance:} self-service needs an explicit route to
    human help when an issue cannot be resolved.
\end{enumerate}

The engineering problem is to combine natural spoken interaction with traceable
information and predictable backend decisions. Language generation alone cannot
establish personal eligibility or official approval, and these distinctions guide
the system's tool boundaries.

\section{Project Objectives}
\begin{enumerate}
  \item \textbf{Deliver the voice-to-tool path:} connect browser audio to LiveKit
    and Gemini Live, retrieve scheme context through the API, and return a spoken
    answer.
  \item \textbf{Publish a reusable catalogue:} represent scheme descriptions,
    citations, documents, application steps, and optional rules in one validated
    schema, and publish an initial reference set.
  \item \textbf{Implement deterministic guidance:} evaluate rule conditions and
    conditional document requirements with explicit outcomes.
  \item \textbf{Preserve accountable publication:} retain every version and
    associate publication with a named reviewer and date.
  \item \textbf{Support human assistance:} create queued handoffs and expose
    operator status updates and context.
  \item \textbf{Verify and measure:} provide automated tests, linting, migrations
    and containers, and measure retrieval accuracy, throughput, and live voice
    behaviour.
\end{enumerate}

\subsection{Final Scope and Changes from the Proposal}
Table~\ref{tab:scope} compares the proposal with the delivered system.

\begin{table}[htbp]
\centering\small
\begin{tabularx}{\textwidth}{@{}P{0.21\textwidth}YY@{}}
\toprule
\textbf{Area} & \textbf{Original proposal} & \textbf{Final system} \\
\midrule
Citizen access & Ordinary telephone call through a SIP bridge.
  & Scroll-driven landing page leading to a browser phone-call screen with
    live captions over LiveKit; no telephony. \\
Scheme guidance & Facts from a verified, published catalogue.
  & 18 records published with a named reviewer attestation; qualified general
    guidance is permitted when retrieved context is incomplete. \\
Eligibility & Deterministic pre-checks over reviewed rules.
  & Engine implemented and tested; 27.8\% of published records carry rule sets. \\
Source ingestion & Import, AI extraction, human review, versioned publish.
  & Structured JSON preparation, reviewer attestation, and versioned
    publication; a full source importer remains future work. \\
Delivery and scale & Staging pipeline and measured replica comparison.
  & CI checks, containers, and migrations; 1-, 2- and 4-replica throughput
    measured locally; staging CD remains future work. \\
\bottomrule
\end{tabularx}
\caption{Relationship between the proposal and the final system.}
\label{tab:scope}
\end{table}

% ================================================================
\chapter{Methodology and Design}

\section{System Architecture}
Janmitra uses a modular backend with a separate voice worker. The backend contains
conversation, catalogue, eligibility, handoff, and audit modules. The voice worker
calls the backend through typed HTTP tools and never writes to the database
directly or duplicates eligibility logic \cite{architecture}.

\begin{figure}[htbp]
\centering
\setlength{\fboxsep}{8pt}
\fbox{\parbox{0.83\textwidth}{\centering
  \textbf{Browser landing page and phone-call screen}\\
  Scheme rings, microphone, speaker, call controls, and live captions}}
\par\smallskip $\updownarrow$\quad WebRTC audio\par\smallskip
\fbox{\parbox{0.83\textwidth}{\centering
  \textbf{LiveKit room and Python voice worker}\\
  Gemini Live native-audio session; tool invocation; transcript and timing events}}
\par\smallskip $\updownarrow$\quad Authenticated HTTP tools\par\smallskip
\fbox{\parbox{0.83\textwidth}{\centering
  \textbf{FastAPI backend}\\
  Conversation state; catalogue retrieval; deterministic eligibility;\\
  document checklists; handoff queue; publication; audit}}
\par\smallskip $\updownarrow$\quad Structured data access\par\smallskip
\fbox{\parbox{0.83\textwidth}{\centering
  \textbf{PostgreSQL}\\
  Published scheme versions, conversations, events, handoffs, and audit records}}
\caption{Final system architecture. Admin and operator clients access the backend
separately.}
\label{fig:architecture}
\end{figure}

\section{Technical Stack}
\begin{itemize}
  \item \textbf{Backend:} Python 3.12 or later, FastAPI, and Uvicorn.
  \item \textbf{Validation and persistence:} Pydantic, asynchronous SQLAlchemy,
    PostgreSQL 16 with asyncpg, and Alembic migrations. SQLite isolates unit tests.
  \item \textbf{Voice:} LiveKit Agents and its Google plugin at version 1.8.5.
    The native-audio model is \code{gemini-3.1-flash-live-preview}.
  \item \textbf{Frontend:} Next.js 16, React 19, TypeScript, and the LiveKit
    client components.
  \item \textbf{Tooling:} pytest, Vitest, Ruff, ESLint, Docker Compose, and
    GitHub Actions.
\end{itemize}

\section{Implementation Details}

\subsection{Voice Conversation and Session Lifecycle}
The worker joins a LiveKit room, creates a backend conversation, and starts an
\code{AgentSession} with Gemini Live. Speech input and output are handled within a
single model session, without a separate speech-to-text or text-to-speech stage.
The citizen speaks first, and the first reply briefly discloses that Janmitra is
an AI. The prompt asks the agent to match the citizen's language, use short spoken
steps, and ask one question at a time.

Voice activity detection uses 100~ms prefix padding and 500~ms silence
endpointing. Transcript and timing events are written asynchronously, and a single
cleanup task drains pending writes and closes the conversation. Status messages
over a data channel report readiness and recoverable errors to the browser.

\subsection{Typed Backend Tools}
The agent uses four operations:
\begin{description}[leftmargin=1.7em,style=nextline]
  \item[\code{find_service}] Retrieves scheme context for the citizen's stated need.
  \item[\code{check_eligibility}] Evaluates answers against a published rule set.
  \item[\code{get_documents}] Returns the published document checklist, including
    conditional items.
  \item[\code{request_handoff}] Queues a request for human assistance.
\end{description}
The backend validates each payload, checks the caller's role, and returns a
structured result. The model explains that result in speech; it cannot publish
data or override the eligibility engine.

\subsection{Scheme Catalogue}
Each \code{ServiceRecord} holds a slug, localised name and description, aliases,
category, benefit and eligibility summaries, an optional rule set, documents,
application steps, and citation metadata. The catalogue contains 18 central
schemes, summarised in Table~\ref{tab:catalogue}.

\begin{table}[htbp]
\centering\small
\begin{tabularx}{\textwidth}{@{}lrY@{}}
\toprule
\textbf{Category} & \textbf{Records} & \textbf{Schemes} \\
\midrule
Loan & 4 & KCC, PMMY, PMEGP, PM-Vishwakarma \\
Insurance & 4 & PM-JAY, PMFBY, PMJJBY, PMSBY \\
Grant & 4 & PM-KISAN, PMAY-G, PMGKAY, PMUY \\
Pension & 3 & APY, PM-SYM, NSAP \\
Banking & 2 & PMJDY, DAY-NRLM \\
Other & 1 & Rural employment guarantee (slug \code{mgnrega}) \\
\midrule
\textbf{Total} & \textbf{18} & All published \\
\bottomrule
\end{tabularx}
\caption{Published scheme catalogue by category.}
\label{tab:catalogue}
\end{table}

\subsection{Retrieval and Grounding}
Retrieval combines Unicode-aware normalisation, Hindi and Punjabi phrase
expansion, transliteration, named-scheme matching, and a curated mapping from
expressed needs to candidate schemes. Compound requests keep each distinct need
rather than returning variants of one topic. When no need is recognised, a
complete alias or scheme name must appear in the query; this deliberately prevents
incidental word overlap in a personal story from being treated as a scheme match.

The voice path requests \code{response_mode="context"}, and the API returns at most
three records with citations and version identifiers. Gemini Live then answers
directly from this context, avoiding an intermediate text-generation round trip.

\subsection{Deterministic Eligibility and Documents}
Rules are stored as data within a versioned record. Named conditions combine
through \code{all_of}, \code{any_of}, and \code{none_of} groups, and optional
\code{applies_when} guards control relevance. The engine validates answers and
returns \code{eligible}, \code{not_eligible}, or \code{needs_more_info}, with a
trace of the conditions involved. Document checklists use the same answers: a
document whose condition is unknown remains visible as conditional.

\subsection{Versioned Publication and Review}
Publishing appends a new \code{ServiceVersion} and moves the service's current
pointer, retaining every earlier version with field-level diffs. The seed command
validates each record and publishes only when a reviewer supplies a name with
\code{--actor} and attests to review with \code{--confirm-reviewed}. PostgreSQL
advisory and row locks serialise concurrent publication. The 18 records were
published through this path, with the reviewer recorded in the version history
\cite{catalogue-review}.

\subsection{Human Handoff and Operator Queue}
A handoff stores an issue summary, optional contact details, language, and the
conversation reference, then marks the conversation \code{handed_off}. Operators
list queued requests, inspect conversation events, add notes, and move requests
from \code{new} to \code{contacted} to \code{resolved}. A handoff queues operator
attention; it does not transfer a live call.

\subsection{Persistence, Audit, and Access Control}
The database holds services, versions, source snapshots, conversations,
conversation events, handoff requests, and audit events. Audit events record
critical actions with request identifiers and actors. Separate admin, operator,
and voice API keys restrict access by role. The prompt forbids requesting Aadhaar
numbers, full bank details, OTPs, PINs, and passwords.

\subsection{Frontend: Landing Page and Call Screen}
The landing page opens on five concentric rings, one per scheme category, each
carrying the names of its published schemes around the Janmitra wordmark. Scrolling
flies the viewer into the rings, with outer rings opening faster as in a tunnel,
until the innermost ring frames three statements in turn: the hardship people
face, why existing help fails to reach them, and how Janmitra answers in the
citizen's own language. The page then ends on the
call screen. Without JavaScript, or when the visitor prefers reduced motion, the
same content appears as ordinary stacked sections.

Scheme names are placed at angles solved against their measured sizes, so that each
stays inside the screen and clear of the wordmark, the call action, and every other
name. A runtime rule hides any name that would still be cut off on an unusual screen
size. The page's catalogue data, scroll timeline, and placement rule live in a pure
module, which makes them testable without a browser.

The call screen reproduces an iPhone in-call screen: a notch and status bar, the
caller name and a live status (listening, thinking, speaking), a grid of mute,
speaker, caption and volume controls, and a single call button. Live captions show
the most recent turns of both speakers. Readiness and error messages from the
worker appear in an accessible alert region.

\subsection{Development and Deployment}
GitHub Actions run Ruff and pytest with a PostgreSQL 16 service, and separately run
frontend linting, tests, a production build, and a dependency security audit. Docker Compose defines PostgreSQL
and the API with health checks; the API container applies migrations before
starting.

% ================================================================
\chapter{Results and Evaluation}

\section{Testing Methodology}
Automated tests were run on \ReportDate\ at revision \texttt{\Revision}, locally
and in continuous integration, with dependencies installed at the versions the
project declares. Retrieval, load, and voice measurements were taken on the same
backend code at earlier revisions; the backend has since changed only in a test
configuration setting and in citation links
\cite{implementation}. Table~\ref{tab:tests} lists the tests performed.

\begin{table}[htbp]
\centering\small
\begin{tabularx}{\textwidth}{@{}P{0.27\textwidth}YP{0.17\textwidth}@{}}
\toprule
\textbf{Test} & \textbf{What it verifies} & \textbf{Tool} \\
\midrule
Backend unit tests & Rule logic, numeric validation, conditional documents,
  ranking, schema constraints, review attestation, handoff triggers & pytest \\
API integration tests & Role restrictions, retrieval, eligibility, documents,
  handoffs, conversation events, audit & pytest, SQLite \\
PostgreSQL integration & Migration round trip, schema agreement, concurrent
  publication, event numbering, competing state transitions & pytest, PostgreSQL 16 \\
Voice-runtime tests & Worker cleanup, notifications, session behaviour, prompt
  and client contracts & pytest \\
Frontend tests & Token route security, voice-event handling, and the landing
  page: every published scheme on its category's ring, the scroll timeline on ten
  screen shapes, and the name placement rule & Vitest \\
Browser layout check & Every scheme name visible, inside the screen, and clear of
  the wordmark, call action, and other names at six screen sizes & Chrome,
  Playwright \\
Continuous integration & Backend, frontend, and documentation workflows on the
  final revision & GitHub Actions \\
Static analysis & Code style and correctness rules & Ruff, ESLint \\
Production build & Type checking and compilation & Next.js \\
End-to-end API checks & Health, authentication, catalogue, conversation
  lifecycle, all four tools, handoff queue, audit trail & Scripted HTTP client \\
Catalogue coverage & Structural completeness of published records & Script \\
Source link check & Reachability of every cited official source & HTTP client \\
Retrieval accuracy & Recall at three results over labelled citizen queries in
  English, Hindi, Punjabi, and transliteration & Scripted benchmark \\
Load and scaling & Throughput, latency, and errors at 1, 2, and 4 replicas &
  Threaded load generator \\
Live voice trials & Complete browser calls in English and Punjabi & Browser,
  LiveKit, Gemini Live \\
Direct model probes & Gemini Live with the worker's configuration, streamed
  audio, and a tool call, without LiveKit & google-genai SDK \\
\bottomrule
\end{tabularx}
\caption{Tests performed for the final evaluation.}
\label{tab:tests}
\end{table}

\section{Summary of Results}
Table~\ref{tab:results} summarises the outcome of each test.

\begingroup
\small
\renewcommand{\arraystretch}{1.12}
\begin{longtable}{@{}P{0.36\textwidth}P{0.18\textwidth}P{0.38\textwidth}@{}}
\caption{Final evaluation results.}
\label{tab:results}\\
\toprule
\textbf{Measure} & \textbf{Result} & \textbf{Note} \\
\midrule
\endfirsthead
\toprule
\textbf{Measure} & \textbf{Result} & \textbf{Note} \\
\midrule
\endhead
\bottomrule
\endfoot
Backend test pass rate & 100\% & Includes PostgreSQL and voice-runtime modules;
  none excluded \\
Frontend test pass rate & 100\% & Includes the landing page tests \\
Scheme names unobstructed & 100\% & At every tested screen size \\
CI workflows passing & 100\% & Backend, frontend, and documentation \\
Static analysis & Passed & Ruff and ESLint report no issues \\
Production build & Passed & \\
End-to-end API checks & 96.9\% & The one unmet check is a housing phrasing
  discussed in Section~\ref{sec:retrieval} \\
Records published & 100\% & With a named reviewer attestation \\
Records with citations & 100\% & \\
Cited source links reachable & 100\% & \\
Records with application steps & 100\% & \\
Records with document checklists & 83.3\% & \\
Records with eligibility rule sets & 27.8\% & APY, PM-SYM, PMJJBY, PMSBY, PMUY \\
Retrieval recall at three results & 97.1\% & Median lookup 4.95~ms \\
Throughput targets met & 100\% & More than tenfold margin at every replica count \\
Load-test error rate & 0\% & \\
Live voice calls completed & 50\% & Remainder closed by provider capacity limit \\
Retrieval success within live calls & 100\% & Every lookup under 50~ms \\
Direct model probes completed & 100\% & \\
\end{longtable}
\endgroup

\section{Retrieval Accuracy}
\label{sec:retrieval}
Retrieval was measured against labelled queries phrased as a citizen would speak
them, with the acceptable schemes fixed in advance. A query passes when a correct
scheme appears among the first three results. Table~\ref{tab:retrieval} gives the
results by query group.

\begin{table}[htbp]
\centering\small
\begin{tabularx}{\textwidth}{@{}lrY@{}}
\toprule
\textbf{Query group} & \textbf{Recall@3} & \textbf{Example} \\
\midrule
English, natural phrasing & 100\% & ``My crop was destroyed by heavy rain'' $\rightarrow$ PMFBY \\
Hindi (Devanagari) & 100\% & A ration request in Devanagari script $\rightarrow$ PMGKAY \\
Punjabi (Gurmukhi) & 100\% & A request for housing $\rightarrow$ PMAY-G \\
Transliterated Hinglish & 100\% & ``mujhe ghar chahiye'' $\rightarrow$ PMAY-G \\
Scheme named directly & 83.3\% & ``What is MUDRA loan'' $\rightarrow$ PMMY \\
\midrule
\textbf{Overall} & \textbf{97.1\%} & Median latency 4.95~ms \\
\bottomrule
\end{tabularx}
\caption{Retrieval recall at three results by query group.}
\label{tab:retrieval}
\end{table}

Every multilingual group reached 100\%, which provides measured evidence for
retrieval across Hindi, Punjabi, and transliterated input. The unmatched queries
follow from the strict alias rule described in Chapter~2: ``Tell me about Ujjwala''
does not match the stored alias ``Ujjwala 2.0'', and a housing need phrased as
storm damage to a roof is not yet in the housing vocabulary, while ``I need a pucca
house'' is matched correctly. The system therefore prefers returning nothing to
returning an unrelated scheme. Extending aliases and need phrases addresses these
cases without weakening that guard.

\section{Performance and Scalability}
Four API replicas were run against a shared PostgreSQL instance with the mock model
adapter, so no provider calls were involved. A load generator issued retrieval
requests in round robin across the active replicas, with a discarded warm-up before
each phase. Table~\ref{tab:throughput} compares the results with the targets set in
the proposal.

\begin{table}[htbp]
\centering\small
\begin{tabular}{@{}rrrrrr@{}}
\toprule
\textbf{Replicas} & \textbf{Target (req/s)} & \textbf{Measured (req/s)} &
  \textbf{\% of target} & \textbf{p50 (ms)} & \textbf{p95 (ms)} \\
\midrule
1 & 20 & 244.7 & 1,224\% & 130.5 & 153.0 \\
2 & 35 & 387.6 & 1,107\% & 83.5 & 162.3 \\
4 & 55 & 559.8 & 1,018\% & 29.8 & 175.4 \\
\bottomrule
\end{tabular}
\caption{Throughput and latency by replica count. Error rate was 0\% in every phase.}
\label{tab:throughput}
\end{table}

Every target was exceeded by more than tenfold with no failed requests. Scaling
efficiency was 79\% at two replicas and 57\% at four. Median latency fell sharply
as replicas were added while aggregate throughput rose less, which indicates that
the single-machine load generator, rather than the API, became the limit. The
measured values are therefore lower bounds on the system's capacity.

\section{Live Voice Evaluation}
Live trials used complete browser calls through LiveKit and the real Gemini Live
model, in English and Punjabi. Backend retrieval succeeded in every call (100\%),
with each lookup completing in under 50~ms, and worker startup had a median of
63.5~ms \cite{voice-design}. Half of the integrated calls (50\%) completed; in the others, the provider
closed the connection with WebSocket code 1011 and the message ``Resource has been
exhausted''.

These failures were traced to the provider's free-tier capacity rather than to the
application. The worker's exact configuration (model, thinking level,
temperature, voice, and voice activity settings) was replayed outside LiveKit with
streamed 16~kHz speech and the real retrieval payload injected mid-session, and
every probe with properly endpointed speech completed (100\%). A text request on
the same key also succeeded. One failure began 24 seconds after a two-minute call
ended, which is consistent with a rate limit that recovers over time. Google does not publish fixed
Live API free-tier limits, so the final demonstration should use a billed project.

In the isolated probe, the first spoken audio arrived about 3.1~s after the end of
the citizen's speech, including a retrieval tool call. This single measurement
falls within the proposal's seven-second time-to-guidance target, but it is not a
median over integrated calls.

The Punjabi call provided the strongest qualitative result. The system transcribed
Gurmukhi speech correctly, translated the need into English for retrieval, matched
PM-KISAN, KCC, and the rural employment scheme, and replied in Punjabi with an
appropriate greeting. In one English call, voice activity detection split the
opening sentence at a short pause, and the reply was generated in Hindi.

\section{Analysis}

\subsection{Achievement of Objectives}
\begin{table}[htbp]
\centering\small
\begin{tabularx}{\textwidth}{@{}P{0.32\textwidth}P{0.15\textwidth}Y@{}}
\toprule
\textbf{Objective} & \textbf{Status} & \textbf{Evidence} \\
\midrule
Voice-to-tool path & Achieved & Retrieval succeeded in 100\% of live calls;
  completed calls returned spoken guidance \\
Reusable, published catalogue & Achieved & 100\% of records published with
  citations and steps \\
Deterministic guidance & Achieved & Engine and checklist tests pass; 27.8\% of
  records carry rule sets \\
Accountable publication & Achieved & Reviewer and date recorded; versions
  retained \\
Human assistance & Achieved & Handoff creation and operator listing verified
  end to end \\
Verification and measurement & Achieved & 100\% test pass rate; retrieval,
  load, and voice measured \\
Ordinary-phone access & Not achieved & Browser channel only \\
\bottomrule
\end{tabularx}
\caption{Status of project objectives.}
\label{tab:objectives}
\end{table}

\subsection{Challenges and Design Responses}
\textbf{Balancing helpfulness and verification.} A verified-only path cannot answer
until data is reviewed. Publishing the catalogue under a reviewer attestation
enables verified tools, while qualified general guidance covers gaps without
inventing amounts, deadlines, or approvals.

\textbf{Understanding ordinary and mixed needs.} Citizens describe situations rather
than scheme names. Need concepts, multilingual expansion, and combined ranking gave
97.1\% recall, and the strict alias rule keeps unrelated matches out.

\textbf{Voice latency.} Answering directly from retrieved context removes a
text-generation round trip, and retrieval stayed under 50~ms in every live call.

\textbf{Consistency under concurrency.} Database locks protect publication,
event numbering, and terminal state transitions, and the PostgreSQL integration
tests verifying them pass at 100\%.

\subsection{Limitations}
\begin{enumerate}
  \item The browser channel does not provide ordinary-phone access, and handoff
    does not transfer a live call.
  \item Free-tier provider capacity limited live call completion to 50\%.
  \item Turn-level timing fields used by the worker are not emitted by the SDK for
    native-audio models, so a median turn-level time-to-guidance could not be
    computed from integrated calls.
  \item Language matching is not yet reliable when voice activity detection splits
    an utterance; the English-language rule is present in the prompt but placed late.
  \item Document checklists cover 83.3\% of records and rule sets 27.8\%.
  \item Handoff precision and recall, claim-level citation accuracy, and user
    clarity require labelled data, source review, and a participant pilot, and were
    not evaluated.
\end{enumerate}

% ================================================================
\chapter{Conclusion}

\section{Summary of Work}
Janmitra has progressed from a proposed phone-accessible guidance system to a
working browser voice service on a tested, modular backend. The delivered system
includes typed guidance tools, a published 18-scheme catalogue, deterministic
eligibility and document logic, versioned publication with reviewer attestation,
an operator handoff queue, audit trails, migrations, containers, CI, a
scroll-driven landing page, and a phone-style call screen with live captions.

The final evaluation found a 100\% pass rate across backend and frontend tests
and every CI workflow,
96.9\% of end-to-end API checks passing, 97.1\% retrieval recall with 100\% on
every multilingual group, and throughput above every scaling target by more than
tenfold at a 0\% error rate. Live voice calls retrieved correctly in every case and
completed in 50\% of trials under free-tier provider limits.

\section{Key Findings}
\begin{itemize}
  \item \textbf{Explicit tool boundaries make guidance testable.} Eligibility,
    publication, and handoff can be verified independently of speech generation.
  \item \textbf{Strict matching trades recall for trust.} Requiring a recognised
    need or a full alias kept unrelated schemes out, at a small cost in recall.
  \item \textbf{The backend is not the bottleneck.} A single replica sustained more
    than twelve times its target, and live lookups completed in under 50~ms.
  \item \textbf{Provider capacity dominates live reliability.} Every integrated
    failure came from the speech provider's limit, while the application path
    completed whenever capacity was available.
\end{itemize}

\section{Future Work and Recommendations}
\begin{enumerate}
  \item \textbf{Telephone access:} add a SIP bridge to meet the original
    ordinary-phone objective.
  \item \textbf{Reliable live capacity:} use a billed Gemini project and lengthen
    the reconnection backoff so retries do not compound capacity limits.
  \item \textbf{Latency instrumentation:} measure time to first audio directly in
    the worker to report a median turn-level time-to-guidance.
  \item \textbf{Language and turn handling:} move the language-matching rule
    earlier in the prompt and tune voice activity endpointing for natural pauses.
  \item \textbf{Catalogue growth:} add state-specific schemes, complete document
    checklists, extend rule coverage, and widen aliases and need phrases.
  \item \textbf{Pilot evaluation:} run the 16 evaluation scenarios with human
    review \cite{evaluation-fixtures}, then a pilot with 8--12 participants
    measuring completion, time, and clarity.
\end{enumerate}

\section{Final Remarks}
Janmitra shows how conversational AI can be connected to structured public-service
information and deterministic backend decisions. Its central engineering
contribution is the separation of speech, retrieval, trusted decisions, and
persistent records into components that can each be tested and measured. With
telephone access and a citizen pilot, the same foundation can move toward the
proposal's goal of accessible and dependable civic guidance.

% All references are embedded so this deliverable remains a single .tex file.
\clearpage
\phantomsection
\addcontentsline{toc}{chapter}{Bibliography}
\begin{thebibliography}{9}

\bibitem{proposal}
Dhruv Srivastava, Harkamal Singh Lubana, and Paras.
\emph{Janmitra: A Voice-First Civic Scheme Guidance Platform for Rural India}.
UCS503P project proposal. Repository source: \path{project-proposal/main.tex}.

\bibitem{prototype}
Dhruv Srivastava, Harkamal Singh Lubana, and Paras.
\emph{Janmitra: Prototype-Stage Project Report}. Mid-semester evaluation,
10 October 2026. \path{project-report-prototype-stage/main.tex}.

\bibitem{implementation}
Janmitra project team. \emph{Implementation and commit history}.
Repository revision \texttt{\Revision}. Backend under \path{code/backend/},
frontend under \path{code/frontend/}, CI in \path{.github/workflows/}.

\bibitem{architecture}
Janmitra project team. \emph{Architecture}. \path{docs/architecture.md}.

\bibitem{voice-design}
Janmitra project team. \emph{Browser Voice Architecture}.
\path{docs/voice-architecture.md}.

\bibitem{catalogue-review}
Janmitra project team. \emph{Catalogue Review Checklist}.
\path{code/backend/data/REVIEW.md} and the scheme JSON records.

\bibitem{evaluation-fixtures}
Janmitra project team. \emph{Voice Evaluation Scenarios and Rubric}.
\path{code/backend/scripts/fixtures/voice-evaluation.json}.

\end{thebibliography}
\end{document}
